\documentclass{article}
\usepackage[utf8]{inputenc}
\usepackage[T1]{fontenc}
\usepackage{geometry}
\geometry{
    a4paper,
    left=2.5cm,
    right=2.5cm,
    top=2.5cm,
    bottom=2.5cm
}
\usepackage{graphicx}
\usepackage{float}
\usepackage{hyperref}
\usepackage[english]{babel}
\usepackage{amsmath}
\usepackage{amssymb}
\usepackage{array}
\usepackage{booktabs}
\usepackage{caption}
\usepackage{subcaption}
\usepackage{enumitem}
\usepackage{fancyhdr}
\usepackage{titlesec}
\usepackage{url}
\usepackage{textcomp}
\usepackage{xcolor}
\usepackage{listings}
\usepackage{multirow}
\usepackage{longtable}
\usepackage{tabularx}
\usepackage{times}

% Set up page style
\pagestyle{fancy}
\fancyhf{}
\rfoot{\thepage}
\renewcommand{\headrulewidth}{0pt}

% Title formatting
\titleformat{\chapter}[display]{\normalfont\bfseries}{\chaptertitle\thechapter}{0pt}{\Huge}
\titlespacing*{\chapter}{0pt}{50pt}{40pt}

% Document info
\title{Template Document}
\author{Academic Support Agent}
\date{\today}

\begin{document}
\maketitle

% Content from preambel would be included here
 	% Alle Packages und Formatierungen werden in preamble.tex eingetragen um den Text hier  bersichtlicher zu gestalten.
 
\section{Einleitung} %  berschrift f r erstes Kapitel, optional auch: [Eintrag im IHV]{Einleitung}
\label{sec:Einleitung} % Textmarke, um darauf zu verweisen.
\lipsum[2]			% Dummytext (einfach l schen und los schreiben)

\section{Thematik in der Literatur}
\lipsum[30]				% Dummytext (einfach l schen und los schreiben)

Es gibt eine Reihe von M glichkeiten chemische Formeln wie in \textsl{\ref{XXX}} zu notieren. Eine davon ist diese:

\begin{equation}					% Gleichung
\begin{aligned}						% Nummerierung am Rand
\label{XXX}							% beliebiges Label
\textnormal{XXX}_2 \hspace{0.3cm}\autorightarrow{dr ber}{drunter}\hspace{0.3cm} \textnormal{XX}_2 + \textnormal{X}
\end{aligned}
\end{equation}
\vspace{0.2cm}

Wenn nur einzelne Ionen notiert werden sollen k nnen diese entweder direkt im Text als $\anion{--II}{O}{}{2}$ geschrieben, oder in der Reaktionsgleichung \textsl{\ref{YYY}} notiert werden:

\begin{equation}					% Gleichung
\begin{aligned}						% Nummerierung am Rand
\label{YYY}							% beliebiges Label
\textnormal{O}_2 + 4\, \ee{}{} \hspace{0.3cm}\autorightarrow{Reduktion}{}\hspace{0.3cm} 2\,\anion{--II}{O}{}{2} 
\end{aligned}
\end{equation}
\vspace{0.2cm}

\begin{figure}[!ht]		% Bild mit Unterschrift einbinden
\centering
\includegraphics[scale=1]{pics/bild.png} 
\caption[Demobild]{Einzelne Bilder wie dieses sollten zentriert auf der Seite platziert werden. Auf eine sinnvolle Gr e ist ebenso zu achten, wie auf eine konkrete Beschreibung der Abbildung.}
\label{fig:bild1}\vspace{0.8em}
\end{figure}

\lipsum[31]

\newpage
\section{ berschrift}

\renewcommand{\baselinestretch}{1.0}\normalsize 
\setlength{\columnsep}{30px} % Abstand Spalten vergr ern
\begin{multicols}{2}[\noindent \textbf{Hier k nnte beispielsweise  ein Verweis auf den Textauszug aus der Literatur stehen.\citeup{voss}{}}] %\textbf{Dies kann als Zusammenfassung f r den darunter liegenden Text gesehen werden. Hin und wieder ist dieser auch fett gedruckt.}
\setcounter{linenumber}{1}
\begin{linenumbers} % Schaltet Zeilennummerierung ein
%\modulolinenumbers[5] % nur jede 5. Zeile
\lipsum[99]
\end{linenumbers}
\end{multicols} % Ende 2 Spalten
\renewcommand{\baselinestretch}{1.5}\normalsize 

\subsection{ berschrift}
\lipsum[99]

\begin{table}[!ht]
\centering
\caption[Zahlenvergleich der Werte 1 und 2\citeup{voss}{}]{Die Werte sind an fiktiven Proben gemessen.}\vspace{0.8em}
\label{tab:tab1}
\begin{tabular}{l|cc} \toprule
 &\multicolumn{2}{r}{universal}\\ \cmidrule{1-1}
\backslashbox{ABC}{DEF} & 1 & 2 \\ \midrule 
E & 3,34 & 4,34 \\
F & 5,45 & 6,45 \\ \bottomrule
\end{tabular}
\end{table}

\lipsum[99]

\newpage
\section{Ausblick und Zusammenfassung}
Dies ist eine Text\footnote{Dies solle Fu note werden, die den Rahmen einer Zeile sprengt. Dies l sst sich auch durch Wiederholung den Inhalts erreichen. Dies l sst sich auch durch Wiederholung den Inhalts erreichen.  Dies l sst sich auch durch Wiederholung den Inhalts erreichen. \newline Dies l sst sich auch durch Wiederholung den Inhalts erreichen.}, der auch in der Literatur existiert.\citeup{demo}{,S.22} 
\lipsum[32]

Da pack ich doch gleich no ne Fu note drunter.\footnote{Mal sehen, welcher Abstand hier zu vorherigen Fu note entsteht.}

\lipsum[10]

\begin{table}[!ht]					%	Tabelle genau hier einsetzen
\centering							% Tabelle zentrieren
\caption[YYY]{Tabellen berschrift.}\vspace{0.8em} % Tabellenbeschrift.
\label{tab:2}							% Label f r den Tabellenverweis
\begin{tabular}{l|r|r} \toprule		% l = linksb., c = zentriert, r = rechtsb., | vert. Strich
  & \textbf{Spalte A} & \textbf{Spalte B} \\ \midrule 				% & ist Spaltentrennzeichen, Abschluss der Zeile mit "\\"
\textbf{Zeile 1} & Wert 1A & Wert 1B \\
\textbf{Zeile 2 }& Wert 2A & Wert 2B \\ \bottomrule		% top-, mid- & bottomrule in untersch. Strichst rken
\end{tabular}
\end{table}

\lipsum[41]


\newpage
\section{Selbst ndigkeitserkl rung}
Hiermit erkl re ich, dass ich die vorliegende Hausarbeit selbst ndig verfasst und keine anderen als die angegebenen Hilfsmittel benutzt habe.\footnote{eine kleine Fu note.}

Die Stellen der Hausarbeit, die anderen Quellen im Wortlaut oder dem Sinn nach entnommen wurden, sind durch Angaben der Herkunft kenntlich gemacht. Dies gilt auch f r Zeichnungen, Skizzen, bildliche Darstellungen sowie f r Quellen aus dem Internet.

\vspace{20mm}
\hrule
\vspace{5mm}
\begin{center}
\textbf{Bad Reichenhall}, \today
\end{center}


\newpage


\newpage


\newpage
\addcontentsline{toc}{section}{Literaturverzeichnis}
\small 						%\documentclass[journal = jacsat]{achemso}
							%\setkeys{acs}{articletitle = true}
 %achemso angewchem 


\newpage

\addcontentsline{toc}{section}{Anhang}
\begin{appendix}
\section*{Anhang}
%\subsection{Anhangsteil 1}
%\lipsum[31]
%\subsection{Anhangsteil 2}
\end{appendix}
\end{document}
\end{document}